% Pre-textuais: dedicatoria, agradecimentos, epigrafe, resumo, abstract

% --- Dedicatoria ---
\begin{dedicatoria}
   \vspace*{\fill}
   \begin{flushright}
   \noindent
   \textit{Aos nossos familiares, que sempre acreditaram\\
   no poder transformador da educa\c{c}\~ao t\'ecnica\\
   e da pesquisa aplicada \`a defesa nacional.}
   \end{flushright}
\end{dedicatoria}

% --- Agradecimentos ---
\begin{agradecimentos}
Agradecemos ao Instituto Militar de Engenharia (IME) pela forma\c{c}\~ao t\'ecnica, cient\'\i fica e humana de excel\^encia, que possibilitou a concep\c{c}\~ao e o desenvolvimento deste trabalho.

Agradecemos especialmente aos nossos orientadores --- Prof. Daniel Rodrigues dos Santos, Prof. Leonardo Assump\c{c}\~ao Moreira, Prof. Bruno Eduardo Madeira e Prof. Gabriel da Cruz Fontenelle --- pelas in\'umeras horas de discuss\~ao t\'ecnica, pelas sugest\~oes metodol\'ogicas e pelo rigor acad\^emico transmitidos ao longo de todo o projeto.

\`A Se\c{c}\~ao de Simula\c{c}\~ao da Academia Militar das Agulhas Negras (AMAN), pela demanda real que deu origem a este trabalho e pela disponibilidade em compartilhar sua experi\^encia operacional com caix\~oes de areia f\'\i sicos.

Aos nossos familiares, pela paci\^encia, pelo apoio incondicional e pela compreens\~ao nos in\'umeros fins de semana dedicados a esta obra.

Aos colegas de turma, com quem partilhamos discuss\~oes t\'ecnicas, c\'odigo, caf\'e e amizade.

Por fim, \`a comunidade de \emph{software livre} --- em especial aos mantenedores do NumPy, OpenCV, Open3D e do template abntex2ime --- sem cuja generosidade este trabalho jamais teria sido vi\'avel.
\end{agradecimentos}

% --- Epigrafe ---
\begin{epigrafe}
    \vspace*{\fill}
    \begin{flushright}
        \textit{``A guerra \'e o reino da incerteza;\\
        tr\^es quartos das coisas em que se baseia\\
        a a\c{c}\~ao na guerra est\~ao envoltos\\
        em uma n\'evoa de incerteza maior ou menor.''\\
        --- Carl von Clausewitz, \emph{Da Guerra}, 1832}
    \end{flushright}
\end{epigrafe}

% --- Resumo (PT) ---
\setlength{\absparsep}{18pt}
\begin{resumo}
\SingleSpacing
Este trabalho apresenta o projeto e a implementa\c{c}\~ao de um sistema de Caix\~ao de Areia com Realidade Aumentada voltado ao treinamento militar em opera\c{c}\~oes no terreno, desenvolvido em atendimento a uma demanda real da Se\c{c}\~ao de Simula\c{c}\~ao da Academia Militar das Agulhas Negras (AMAN). O sistema integra um sensor de profundidade Microsoft Kinect~v2, um projetor convencional e um Modelo Digital de Eleva\c{c}\~ao (MDE) de refer\^encia, fornecendo \emph{feedback} visual em tempo real --- vermelho para excesso de areia, azul para falta e verde para conformidade --- diretamente sobre a superf\'\i cie da areia. A calibra\c{c}\~ao geom\'etrica Kinect~$\to$~Mesa \'e obtida de forma emp\'\i rica e autom\'atica a partir de uma superf\'\i cie plana de refer\^encia (uma tampa apoiada sobre o caix\~ao ou a pr\'opria base vazia): como o campo de vis\~ao do sensor \'e mais largo que a \'area \'util do caix\~ao, o algoritmo RANSAC isola primeiro o maior conjunto de pontos coplanares, descartando moldura, piso e ru\'\i do, e a normal do plano \'e ent\~ao refinada por decomposi\c{c}\~ao em valores singulares (SVD) apenas sobre esses \emph{inliers}; na sequ\^encia, constroem-se bases ortonormais por Gram-Schmidt e transforma\c{c}\~oes af\'\i ns homog\^eneas que levam a cota zero \`a base de madeira do caix\~ao. A proje\c{c}\~ao Mesa~$\to$~Projetor emprega o modelo \emph{pinhole} de Tsai. A renderiza\c{c}\~ao emprega malha discretizada de $30 \times 30$ c\'elulas com agrega\c{c}\~ao espacial para filtragem estat\'\i stica do ru\'\i do do sensor. A arquitetura monol\'\i tica modularizada disp\~oe de cadeia de \emph{fallback} em quatro n\'\i veis (PyKinect2, Open3D, libfreenect2 e simula\c{c}\~ao interativa via \emph{mouse}), assegurando opera\c{c}\~ao mesmo na aus\^encia de hardware. A qualidade do motor matem\'atico \'e assegurada por uma su\'\i te de 88 testes unit\'arios automatizados, todos aprovados, cobrindo todos os contratos alg\'ebricos do sistema e o erro quantitativo do modelo de eleva\c{c}\~ao medido. O sistema foi validado tanto em modo simula\c{c}\~ao quanto em ensaios de campo com o sensor Kinect~v2 real, nos quais foi identificado e corrigido um defeito de calibra\c{c}\~ao que restringia a \'area projetada a uma pequena fra\c{c}\~ao da janela, corre\c{c}\~ao formalizada pela introdu\c{c}\~ao de uma matriz de transla\c{c}\~ao $T_{\text{shift}}$ ao centro l\'ogico da mesa.

\textbf{Palavras-chave}: \imprimirpalavraschave
\end{resumo}

% --- Abstract (EN) ---
\begin{resumo}[Abstract]
\begin{otherlanguage*}{english}
\SingleSpacing
This work presents the design and implementation of an Augmented Reality Sandbox system for military terrain-operations training, developed to meet an actual demand from the Simulation Section of the Agulhas Negras Military Academy (AMAN). The system integrates a Microsoft Kinect v2 depth sensor, a conventional projector, and a reference Digital Elevation Model (DEM), providing real-time visual feedback --- red for excess sand, blue for shortage, and green for conformity --- projected directly onto the sand surface. The Kinect-to-table geometric calibration is obtained empirically and automatically from a flat reference surface (a lid placed over the box or the empty base itself): since the sensor's field of view is wider than the box's usable area, a RANSAC algorithm first isolates the largest set of coplanar points, discarding the frame, floor, and noise, after which the plane normal is refined by singular value decomposition (SVD) over those inliers alone; orthonormal bases are then built via Gram-Schmidt and combined into homogeneous affine transformations that place the zero elevation at the box's wooden base. The table-to-projector mapping employs Tsai's pinhole camera model. Rendering employs a $30 \times 30$ discretized grid with spatial aggregation for statistical noise filtering. The monolithic modular architecture features a four-tier fallback chain (PyKinect2, Open3D, libfreenect2, and interactive mouse-based simulation), ensuring operation even without hardware. The mathematical engine's correctness is backed by an automated suite of 88 unit tests, all passing, covering every algebraic contract of the system and the quantitative error of the measured elevation model. The system was validated both in simulation mode and in field trials with the real Kinect v2 sensor, during which a calibration defect that restricted the projected area to a small fraction of the window was identified and corrected, formalized by the introduction of a translation matrix $T_{\text{shift}}$ to the logical table center.

\vspace{\onelineskip}
\noindent\textbf{Keywords}: \imprimirkeywords
\end{otherlanguage*}
\end{resumo}
